\section{Appendix: Scope and Qualifications}
\label{sec:scope}

{\small
\textbf{Edition and Mass.} The study treats the Mass of the First Sunday of
Advent, Year~B, as the \work{Roman Missal}, Third Edition, and the
\work{Lectionary for Mass} give it for the dioceses of the United States, in
the national calendar alone, for the occurrence of 29~November 2026. Identity,
occurrence and branches are set out in the preceding appendix and recorded in
\texttt{instance/manifest.md} and \texttt{research/context.md}. Nothing in the
study holds for a celebration that replaces this Mass with another.

\textbf{Text control.} The Missal's texts were read in the Latin of the
\work{Missale Romanum}, editio typica tertia (2002), and checked against the
list of changes for the 2008 reprint, which names no page of this Sunday; no
page of the 2008 reprint and no United States altar book was inspected, so the
approved English of the prayers was not read in print. The Lectionary's
boundaries rest on the Latin \work{Ordo lectionum Missae} (1981), no.~2, read
on page images, on the national calendars for 2026 and 2027, and on the
bishops' conference's page of readings for 29~November 2026, read on 26
November 2026. The printed United States Lectionary was not inspected. At
Isaiah 64:2 the page of readings cites the whole verse and its text has the
first member (the wonderful deeds) and nothing for \latin{Descendisti}, while
the \work{Ordo} reads the second member only; the bishops' current
translation of the Bible notes at that verse that it omits words the Hebrew
repeats from 63:19. Whether the Lectionary's own translation has the same
note was not established. The \work{Nova Vulgata} was read in the Holy See's
web publication at Isaiah 63:15--65:5 and at Psalms 25 (24), 80 (79) and
85 (84). No Hebrew or Greek text was read: what the Hebrew has at Psalm
80:16b is reported only through the King James Version, and the Septuagint
clause Jerome cites at Isaiah 63:17 is reported through his Latin. The
Entrance Antiphon's \latin{qui te exspectant} is in neither Latin psalter; a
record of Augustine's commentary on the psalm notes that his psalter read
\latin{exspectant}, but that page could not be re-read, and the source of the
antiphon's wording was not traced. The early history of the solemn blessing
and of the first half of the Preface's thanksgiving was not found in what was
read. The branch rows rest on the United States \work{General Instruction}:
chapter~II as read in an Internet Archive capture of its official page made
on 15~October 2026, and chapters~VII and~VIII in captures made on 13 and
18~June 2026, because the publisher's site refused direct requests on
26~November 2026.

\textbf{English.} All Scripture is printed in the Douay--Rheims version as
revised by Bishop Challoner, a public-domain study translation of the
Clementine Vulgate. It is not the text proclaimed in the celebration; the
controlling English is that of the \work{Lectionary for Mass for Use in the
Dioceses of the United States of America}, and for the antiphons, prayers,
Preface and blessing that of the \work{Roman Missal}, Third Edition, neither of
which is reproduced. The prayers, Preface and blessing are identified by Latin
incipit and described; the descriptions are not translations and may not be
quoted as the texts' words. Where the Lectionary appoints part of a verse, the
rest is printed in brackets. English renderings of Latin passages are the
study's own. Augustine on the psalms is quoted in the Oxford translation of
the Nicene and Post-Nicene Fathers, which numbers the sections of Psalm 84
(85) differently from the Latin: the Latin 84.9 and 84.15 are its §§6 and~12.
That translation abridges §6; the sentences of 84.9 it omits are translated
from the Latin.
In quotations the Latin is normalized only by writing out ligatures.

\textbf{Reception.} Each appointed Scripture passage has at least one patristic
commentary read at an exact place. Jerome on Isaiah and on Matthew was read on
page images of the Migne printing; Aquinas on Isaiah and on First Corinthians
in the electronic Latin of the Corpus Thomisticum; Bede on Mark and Gregory's
homilies in electronic Latin texts of the Migne printing; Augustine's letter
199 in the Latin of augustinus.it; Augustine on the psalms, Chrysostom and
Cyril of Jerusalem in the English of the Nicene and Post-Nicene Fathers,
and Augustine on Psalm 84 also in the Latin of the Migne printing;
Bellarmine in O'Sullivan's abridged English of 1866, the Latin not compared.
No text was collated with a modern critical edition. None of the Fathers
cited comments on this Mass as a whole, and no liturgical commentary on it is
cited. Jerome,
Chrysostom and Gregory comment on parallels to the Gospel (Matthew 24, Luke
12); Bede's comments on Mark 13:32--33 are Jerome's, and are credited to
Jerome. The \work{Catena aurea} on Mark was read and is not cited: what it
adds comes from Theophylact, an ecclesiastical writer, whose four watches as
four ages of life differ from Gregory's three, and from Hilary and a writer
under Jerome's name, who are not relied on. Not opened, among others:
Cassiodorus on Psalm 79; Cyril of Alexandria and Theodoret on Isaiah;
Cornelius a Lapide on First Corinthians.

\textbf{Relations among the texts.} Four relations are stated by the liturgical
books: in Advent the readings are chosen for the season, the First Sunday's
Gospel concerning the Lord's coming at the end of time, the prophet's reading
the Messiah, and the apostle's an exhortation for the season
(\work{Praenotanda} 67, 93); the psalm responds to the first reading
(\work{General Instruction} 61); the Alleluia verse greets the Gospel
(\work{General Instruction} 62); and Advent's double character
(\work{Universal Norms} 39) is the Preface's two comings. The \work{Ordo}'s
three titles single out the plea for God's descent, the waiting for Christ's
revelation and the watch. The Missal's antiphons, prayers, Preface and blessing
serve all three years. Among the authors cited, Jerome links Isaiah
64:3 (Douay 64:4) to 1~Corinthians 2:9, and Isaiah 64:6 (Douay 64:7) to Psalm
80 (79):4; Jerome reads the petition of Isaiah 63:17 also as \latin{converte
nos}, ``turn us'', and Aquinas sets Lamentations 5:21 beside it; Augustine
reads the psalm of the Lord's coming;
Augustine and Aquinas link the last day to each one's death. Every other
connection drawn between one text of the Mass and another is the study's own,
grounded in an author's statement about his own passage where one is cited
and otherwise in the wording of the texts; the agreements of Latin words
(\latin{occurrentes} and \latin{Occurristi}; Jerome's \latin{converte nos}
and the psalm's refrain; \latin{exspectant},
\latin{exspectantibus} and \latin{exspectare}; \latin{vigilantes} and
\latin{vigilate}) are observations about the texts. None of them links Isaiah 63--64
to Mark~13.

\textbf{Passages not used.} At cols.~644 and~647 of his commentary on Isaiah
Jerome applies the trodden sanctuary of 63:18 and the soiled justice of 64:5
(Douay 64:6) to the Jewish people after the Passion. Those sentences are not
used, and nothing in the study rests on them. Nor is Chrysostom's comparison,
in the same homily on First Corinthians, of those who refuse the call.

\textbf{Chronology.} The Scriptural Date and Location appendix prints the
shared chronology corpus's answers unchanged, with their sources and their
own qualifications. Its Passage cell for Isaiah prints the Clementine
numbering, which is the Douay--Rheims numbering given beside the Lectionary's
throughout the study.

\textbf{Rights.} The Douay--Rheims, the Clementine Vulgate, the King James
Version, the Nicene and Post-Nicene Fathers, O'Sullivan's Bellarmine, the
volumes of Migne, Wilson's and Feltoe's editions and the \work{Catholic
Encyclopedia} are in the public domain. The Corpus Thomisticum texts are
quoted in short phrases. The Latin and English of the Missal, the \work{Ordo
lectionum Missae}, the Lectionary, the \work{General Instruction}, the national
calendars, the \work{Nova Vulgata}, the \work{Book of Blessings} and the
\work{New American Bible, Revised Edition} are protected; they are cited, and
only incipits, the \work{Ordo}'s three titles and its Latin response and verse,
and short Latin phrases of the prayers needed by the argument are given.

\textbf{Records.} The texts of the Mass are audited in
\texttt{propers/verified.md} and in the collection's record for the First
Sunday of Advent; the reception search, places read, disagreements and
negative results in \texttt{research/scope.md}; and the reasoning behind the
three interpretations, with the class of every link between texts, in
\texttt{research/interpretations.md}. The research records received an
independent review before the study was written.
\par}

\phantomsection
\section*{References}
% A starred heading sets no mark, so without this the Appendix above would go
% on heading the pages of the bibliography. The heading stays starred because
% check-content-preflight finds the References section by that exact form.
\runninghead{References}
\addcontentsline{toc}{section}{References}
\label{sec:references}

{\small
\subsection*{Liturgical books, Bibles and texts}
\begin{itemize}[itemsep=0pt,topsep=.25em]
\item \work{Missale Romanum}, editio typica tertia (Typis Vaticanis, 2002), \latin{Dominica I Adventus}; \latin{Praefatio I de Adventu}; the solemn blessing \latin{In Adventu}; the \latin{Normae universales de anno liturgico et de calendario} 5, 39--40 and 59--60 and the \latin{Calendarium Romanum generale}, as printed in that Missal.
\item \work{Ordo lectionum Missae}, editio typica altera (Libreria Editrice Vaticana, 1981): \work{Praenotanda} 67, 93 and 123; no.~2, \latin{Dominica I Adventus}, Year~B.
\item \work{Lectionary for Mass for Use in the Dioceses of the United States of America}, second typical edition, vol.~I, no.~2, as witnessed by the United States Conference of Catholic Bishops, \work{Liturgical Calendar for the Dioceses of the United States of America} 2026 and 2027, and its page of readings for 29~November 2026, \url{https://bible.usccb.org/bible/readings/112926.cfm}, read on 26~November 2026.
\item \work{General Instruction of the Roman Missal}, United States edition of 2011 with the emendations of 2021, nos.~48, 51, 61--62, 74, 87 (chapter~II), 364--365 (chapter~VII) and 372, 374 (chapter~VIII), read in Internet Archive captures of the official pages of 15~October, 13~June and 18~June 2026.
\item \work{Book of Blessings} (1989), nos.~1509--1520, Order for the Blessing of an Advent Wreath, in the Catholic Culture transcription, read on 26~November 2026.
\item The Holy Bible, Douay--Rheims version, revised by Bishop Richard Challoner (Project Gutenberg text).
\item \work{Biblia Sacra Vulgatae editionis} (Clementine Vulgate), eBible.org text.
\item \work{Nova Vulgata Bibliorum Sacrorum editio}, as published by the Holy See at vatican.va.
\item The Holy Bible, King James Version (eBible.org text).
\item United States Conference of Catholic Bishops, \work{New American Bible, Revised Edition}, Isaiah 64 with its notes, and the introductions to Isaiah and to the Psalms, read in their public deliveries (described; protected).
\end{itemize}

\subsection*{Fathers, Doctors and saints}
\begin{itemize}[itemsep=0pt,topsep=.25em]
\item St Jerome, \work{Commentarii in Isaiam} XVII, on 63:15--64:9 (Migne, PL 24, Paris, 1865, cols.~641--650).
\item St Jerome, \work{Commentarii in Matthaeum} IV, on 24:36 and 24:42--43 (Migne, PL 26, Paris, 1845, cols.~181--182).
\item St Augustine, \work{Enarrationes in Psalmos} 24, 79 and 84, Nicene and Post-Nicene Fathers, 1st ser., vol.~8, and the Latin of the Migne printing for Psalm 84.
\item St Augustine, \work{Epistula} 199, to Hesychius, Latin text of augustinus.it.
\item St John Chrysostom, \work{Homilies on Matthew} 77, Nicene and Post-Nicene Fathers, 1st ser., vol.~10; \work{Homilies on First Corinthians} 2, vol.~12.
\item St Cyril of Jerusalem, \work{Catechetical Lectures} 15, trans. E.~H. Gifford, Nicene and Post-Nicene Fathers, 2nd ser., vol.~7.
\item St Gregory the Great, \work{Homiliae in Evangelia} 9 and 13, Latin text of the Migne printing.
\item St Bede, \work{In Marci evangelium expositio} IV, on 13:32--37, Latin text of the Migne printing.
\item St Thomas Aquinas, \work{Expositio super Isaiam ad litteram}, c.~63--64 (Corpus Thomisticum text).
\item St Thomas Aquinas, \work{Super I epistolam ad Corinthios} c.~1, lect.~1 (Corpus Thomisticum text).
\item St Robert Bellarmine, \work{A Commentary on the Book of Psalms}, trans. J. O'Sullivan (1866, abridged), in the eCatholic2000 transcription: on Psalms 24, 79 and 84.
\end{itemize}

\subsection*{Sacramentaries}
\begin{itemize}[itemsep=0pt,topsep=.25em]
\item H.~A. Wilson, ed., \work{The Gelasian Sacramentary: Liber Sacramentorum Romanae Ecclesiae} (Oxford, 1894), p.~218, for the Collect.
\item C.~L. Feltoe, ed., \work{Sacramentarium Leonianum} (Cambridge, 1896), pp.~20, 74 and 134, for the Prayer over the Offerings, the Prayer after Communion and the Preface.
\end{itemize}

\subsection*{Chronology sources}
\begin{itemize}[itemsep=0pt,topsep=.25em]
\item \work{The Catholic Encyclopedia} (New York, 1908--1912): Howlett, ``Biblical Chronology'', and Reid, ``Captivities of the Israelites'' (vol.~3, 1908); Aherne, ``Epistles to the Corinthians'', and Corbett, ``King David'' (vol.~4, 1908); Maas, ``Chronology of the Life of Jesus Christ'' and ``Jesus Christ'', and Souvay, ``Isaias'' (vol.~8, 1910); MacRory, ``Gospel of Saint Mark'' (vol.~9, 1910); Prat, ``St.~Paul'' (vol.~11, 1911); Drum, ``Psalms'' (vol.~12, 1911); Durand, ``The New Testament'' (vol.~14, 1912); read in the New Advent text.
\item Pontifical Biblical Commission, \latin{De auctore, de tempore compositionis et de historica veritate Evangeliorum secundum Marcum et secundum Lucam}, \work{Acta Apostolicae Sedis} 4 (1912), pp.~463--465, responses I--IX, as the chronology corpus records them.
\end{itemize}
\par}
